% 90-120 minute graduate lecture: camera geometry, two-view geometry and stereo, monocular depth, camera pose and SfM, SLAM, and feed-forward 3D reconstruction.
% Narrative thread: the projection model comes first because every later method either inverts it (depth, stereo, pose) or optimises it over many views (SfM, SLAM, feed-forward models).
\documentclass[10pt,aspectratio=169]{beamer}

\usepackage{tikz}
\usetikzlibrary{positioning, arrows.meta, calc, shapes.geometric, fit, backgrounds, decorations.pathreplacing, patterns}

\input{preamble/packages}
\input{preamble/commands}
\input{preamble/beamer_settings}

\title[Geometric and 3D Vision]{Geometric and 3D Vision}
\subtitle{Depth Estimation, Stereo, Camera Pose, and SLAM}

\begin{document}

\input{sections/cover}
\input{sections/3d-vision/toc}
\input{sections/3d-vision/recap_bridge}
\input{sections/3d-vision/learning_outcomes}
\input{sections/3d-vision/introduction}
\input{sections/3d-vision/camera_geometry}
\input{sections/3d-vision/two_view_geometry}
\input{sections/3d-vision/stereo}
\input{sections/3d-vision/monocular_depth}
\input{sections/3d-vision/camera_pose_sfm}
\input{sections/3d-vision/slam}
\input{sections/3d-vision/feedforward_3d}
\input{sections/3d-vision/limitations}
\input{sections/3d-vision/summary}
\input{sections/3d-vision/references}

\end{document}
